\section{组合与执行纪律}

\begin{exhibitbox}[Exhibit 15：行为框架]
\begin{tabularx}{\textwidth}{L{2.4cm}L{3.3cm}L{3.2cm}X}
\toprule
\textbf{类别} & \textbf{进入条件} & \textbf{加仓／升级} & \textbf{退出／降级} \\
\midrule
正式核心 & 概率目标高于现价且模型可复算 & 下季门槛兑现、回撤承接稳定 & 触发失效条件或概率价值跌破现价 \\
低位盈利 & 一年位置低、扣非为正、增长为正 & 资金由单日变连续、现金流与新研报确认 & Q2恶化、一次性上升、资金继续缺席 \\
已启动候选 & 价格启动、行业阶段确认、盈利桥有效 & H2利润、现金流、订单继续兑现 & 追高后估值超阈值或旧分母失效 \\
条件观察 & 产业方向对但证据或赔率不足 & 补齐ASP／客户／分部毛利／现金流 & 只剩主题热度、牛市情景成为唯一依据 \\
\bottomrule
\end{tabularx}
\end{exhibitbox}

\section{关键风险}

\begin{itemize}
  \item \textbf{预告风险}：业绩预告未经审计，H1不能机械年化；正式中报可能改变扣非、现金流与分部口径。
  \item \textbf{周期风险}：有色、石化、存储和煤炭的高利润可能来自价格、库存或价差，倍数会随周期迅速压缩。
  \item \textbf{资金口径风险}：公开“主力资金”无法穿透最终受益人，静默标签只是价格—资金结构描述。
  \item \textbf{Street时效风险}：徐工、沪农商行和工业富联的正权重原始目标存在时效折价；最新转载页只作零权重核对。
  \item \textbf{高增长边界风险}：恒瑞单品PoS、工业富联客户分配、江波龙合同价、中兴分部毛利属于已检查后的正式披露边界；模型分别以SOTP折价、合并PE、周期PE和条件观察处理。
  \item \textbf{数据传输风险}：申万分类文件下载使用证书校验降级；浙江东方研报接口失败后已回退至官方预告、Q1财务和公开检索，Street权重为0。
\end{itemize}

\section{监控触发器}

\begin{exhibitbox}[Exhibit 16：未来一个季度监控]
\begin{tabularx}{\textwidth}{L{2.1cm}L{3.5cm}L{3.4cm}X}
\toprule
\textbf{对象} & \textbf{升级触发} & \textbf{降级触发} & \textbf{对应动作} \\
\midrule
低位有色／石化 & 现金流转正、价格与资金同步 & 价差回落、库存或现金流恶化 & 从观察升级为分批配置，或退出 \\
非银金融 & 新研报分母、成交与资金持续确认 & 只有低位、没有业务或资金验证 & 保持小权重观察 \\
恒瑞医药 & 创新药增长30\%+、里程碑按期 & 增长低于20\%、EPS低于1.25元 & 回撤核心或下修SOTP \\
工业富联 & H2净利321.1亿元+、毛利率7\%+ & 平台延迟、净利低于516亿元 & 回撤观察或降级 \\
江波龙 & H2现金流转正、合同价与订单延续 & 库存减值、价格回落、现金流继续负 & 重新评估周期情景 \\
中兴通讯 & H2利润正增、算力分部证据出现 & 现金流和利润继续弱 & 维持条件观察 \\
\bottomrule
\end{tabularx}
\end{exhibitbox}

\section{最终配置判断}

当前最值得进入正式决策的是恒瑞医药；渝农商行与徐工机械适合在回撤中复核，工业富联适合等H2兑现，沪农商行偏收益型。全市场低位优先池中，宏桥控股、天山铝业、恒力石化、中孚实业和中矿资源具有明显位置与盈利弹性，但仍需用现金流和行业资金确认。右侧候选中巨人网络质量最好，恒逸石化、东方盛虹与高德红外分别受现金流、估值或旧分母约束。

这不是“所有高增长都买”，而是把全市场机会拆成：正式估值通过、低位盈利待资金确认、已启动待H2兑现、证据已闭环但赔率／现金流未过关四种行为。
